\section{Results and Analysis}
\label{sec:results}

\paragraph{Execution summary and reporting contract.}

All 32 cells trained from initialization on eight \HevHardware{} GPUs in
FP32 (code \texttt{\HevCodeSHAShort}) and used \HevGPUHours{} GPU-hours.
Table~\ref{tab:dev-matrix} reports corrected
development-selection F1. Each cell selects its checkpoint and operating
threshold on the same eight-scene subset. These scores therefore include
model-selection optimism and are not held-out estimates.
Figure~\ref{fig:label-efficiency} shows the four-point curves, and
Fig.~\ref{fig:training-dynamics} shows their training histories in the
appendix. The evidence generator reconstructs every value from the frozen
cohort, checkpoint-bound completion records, and full training curves. Two
cells ran their last development evaluation in a zero-step resume that did
not save the curve; their markers bind the recovery record instead
(Appendix~\ref{app:audit}). The generator rejects an incomplete cohort or a
value that disagrees with its source record.

\begin{table}[t]
\centering
\caption{Development-selection F1 for all 32 cells at each cell's
checkpoint-bound operating threshold. Budgets are fractions of the 111
training scenes (scene counts in parentheses). Per-cell thresholds, best
epochs, and held-out slots appear in Appendix~\ref{app:inventory}.}
\label{tab:dev-matrix}
\small
\setlength{\tabcolsep}{5pt}
\begin{tabular}{@{}llcccc@{}}
\toprule
Track & Initialization & 10\% (12) & 25\% (28) & 50\% (56) & 100\% (111) \\
\midrule
ViT-B/16 & Random & \HevDevFViTRandomTen & \HevDevFViTRandomTwentyFive & \HevDevFViTRandomFifty & \HevDevFViTRandomHundred \\
 & Optical & \HevDevFViTOpticalTen & \HevDevFViTOpticalTwentyFive & \HevDevFViTOpticalFifty & \HevDevFViTOpticalHundred \\
 & SAR & \HevDevFViTSarTen & \HevDevFViTSarTwentyFive & \HevDevFViTSarFifty & \HevDevFViTSarHundred \\
 & ImageNet & \HevDevFViTImageNetTen & \HevDevFViTImageNetTwentyFive & \HevDevFViTImageNetFifty & \HevDevFViTImageNetHundred \\
\midrule
ConvNeXt-V2-B & Random & \HevDevFCNNRandomTen & \HevDevFCNNRandomTwentyFive & \HevDevFCNNRandomFifty & \HevDevFCNNRandomHundred \\
 & Optical & \HevDevFCNNOpticalTen & \HevDevFCNNOpticalTwentyFive & \HevDevFCNNOpticalFifty & \HevDevFCNNOpticalHundred \\
 & SAR & \HevDevFCNNSarTen & \HevDevFCNNSarTwentyFive & \HevDevFCNNSarFifty & \HevDevFCNNSarHundred \\
 & ImageNet & \HevDevFCNNImageNetTen & \HevDevFCNNImageNetTwentyFive & \HevDevFCNNImageNetFifty & \HevDevFCNNImageNetHundred \\
\bottomrule
\end{tabular}
\end{table}

\begin{figure}[t]
\centering
\includegraphics[width=\textwidth]{heldout_label_efficiency.pdf}
\caption{Development-F1 label-efficiency curves, one panel per track. Solid
lines are pretrained arms. The dashed gray line is the track's random floor.
Values are seed-0 point estimates. We draw no uncertainty ribbon.}
\label{fig:label-efficiency}
\end{figure}

\paragraph{Label-efficiency curves and transfer gains.}

At 12 scenes, every pretrained arm beats its random floor in both tracks.
The ViT gains are \HevDeltaFTenViTOptical{} for optical,
\HevDeltaFTenViTSar{} for SAR, and \HevDeltaFTenViTImageNet{} for ImageNet
(Eq.~\eqref{eq:gain} and Fig.~\ref{fig:transfer-gains}). The CNN gains are
\HevDeltaFTenCNNOptical{}, \HevDeltaFTenCNNSar{}, and
\HevDeltaFTenCNNImageNet. CNN-optical now beats random at every budget
because of the input adapter (Sect.~\ref{sec:methods}).

At 12 scenes, ViT-ImageNet comes within 0.002 of the full-data ViT random
floor (\HevDevFViTImageNetTen{} versus \HevDevFViTRandomHundred).

\paragraph{Source-domain contrasts and architecture generality.}

The SAR--optical ranking depends on the backbone
(Eq.~\eqref{eq:contrast} and Fig.~\ref{fig:transfer-gains}). On
development F1, CNN-SAR leads CNN-optical at every budget, but the gap
shrinks from \HevDevFCNNSarTen{} versus \HevDevFCNNOpticalTen{} at 12 scenes
to \HevDevFCNNSarHundred{} versus \HevDevFCNNOpticalHundred{} at 111. The
BigEarthNet Sentinel-1 and Sentinel-2 checkpoints share their architecture,
dataset family, and supervised task, which makes this the cleanest modality
comparison in the study. In the ViT track, optical leads SAR at 10\%, 25\%,
and 50\%, and SAR leads only at 100\%. The scarce-label SAR--optical order
therefore does not replicate across architectures. The ViT checkpoints also
differ in objective and source corpus, so their ranking cannot identify a
modality-only effect.

\paragraph{Held-out test results.}

\ifHevTestComplete
All 32 cells were scored once on the frozen 16-scene test split with their
development-bound thresholds and \HevTestInferencePrecision{} inference.
The test column appears only after every result validates against the frozen
cohort. Test F1 is lower than development F1 for every cell, by 0.029 to
0.142 and \HevTestMeanGap{} on average, as selection on eight scenes leads
one to expect.

Every pretrained arm still beats its random floor at 12 scenes. In the CNN
track, SAR leads optical at 12, 28, and 56 scenes (\HevTestFCNNSarTen{}
versus \HevTestFCNNOpticalTen{} at 12) and trails it at 111
(\HevTestFCNNSarHundred{} versus \HevTestFCNNOpticalHundred). In the ViT
track, the two arms tie at 12 scenes (\HevTestFViTSarTen{} versus
\HevTestFViTOpticalTen), optical leads at 28, and SAR leads at 56 and 111.
The full-data winner changes in one track: ViT-SAR reaches
\HevTestFViTSarHundred, above ImageNet's \HevTestFViTImageNetHundred, while
CNN-ImageNet keeps its lead at \HevTestFCNNImageNetHundred. On test, 28
pretrained scenes already match the full-data random floor in both tracks.
The 28 cells outside the CNN-optical arm repeat the recipe and seed of an
earlier cohort (code \texttt{1a82d508}), yet their test F1 differs from it
by up to 0.043, so we treat contrasts below about 0.04 as run-to-run
variation.

The predeclared monotonicity gate failed on two SAR test curves. ViT-SAR
falls from \HevTestFViTSarTen{} at 10\% to \HevTestFViTSarTwentyFive{} at
25\%, and CNN-SAR falls from \HevTestFCNNSarFifty{} at 50\% to
\HevTestFCNNSarHundred{} at 100\% although its development F1 rises. Both
declines exceed the allowed 0.02, so the experiment does not establish a
monotone label-efficiency curve. The final evaluation proceeded with this
failure disclosed, with no retraining and no threshold tuning.
\else
The 16-scene test split has not been scored in this build. The held-out
contract first freezes the 32-cell cohort. It then applies each cell's
development-bound threshold and writes one immutable result. The generator
admits test values only when all 32 results validate.
\fi

\paragraph{Once-only human-verified evaluation.}

\ifHevFinalEval
\begin{table}[t]
\centering
\caption{F1 on the \HevFinalScenes{} human-verified scenes, scored once
with each cell's development-bound threshold. Appendix~\ref{app:inventory}
adds recall, dark-vessel recall, and near-shore F1.}
\label{tab:final}
\small
\setlength{\tabcolsep}{5pt}
\begin{tabular}{@{}llcccc@{}}
\toprule
Track & Initialization & 10\% (12) & 25\% (28) & 50\% (56) & 100\% (111) \\
\midrule
ViT-B/16 & Random & \HevFinalFViTRandomTen & \HevFinalFViTRandomTwentyFive & \HevFinalFViTRandomFifty & \HevFinalFViTRandomHundred \\
 & Optical & \HevFinalFViTOpticalTen & \HevFinalFViTOpticalTwentyFive & \HevFinalFViTOpticalFifty & \HevFinalFViTOpticalHundred \\
 & SAR & \HevFinalFViTSarTen & \HevFinalFViTSarTwentyFive & \HevFinalFViTSarFifty & \HevFinalFViTSarHundred \\
 & ImageNet & \HevFinalFViTImageNetTen & \HevFinalFViTImageNetTwentyFive & \HevFinalFViTImageNetFifty & \HevFinalFViTImageNetHundred \\
\midrule
ConvNeXt-V2-B & Random & \HevFinalFCNNRandomTen & \HevFinalFCNNRandomTwentyFive & \HevFinalFCNNRandomFifty & \HevFinalFCNNRandomHundred \\
 & Optical & \HevFinalFCNNOpticalTen & \HevFinalFCNNOpticalTwentyFive & \HevFinalFCNNOpticalFifty & \HevFinalFCNNOpticalHundred \\
 & SAR & \HevFinalFCNNSarTen & \HevFinalFCNNSarTwentyFive & \HevFinalFCNNSarFifty & \HevFinalFCNNSarHundred \\
 & ImageNet & \HevFinalFCNNImageNetTen & \HevFinalFCNNImageNetTwentyFive & \HevFinalFCNNImageNetFifty & \HevFinalFCNNImageNetHundred \\
\bottomrule
\end{tabular}
\end{table}

All 32 cells were scored once on the \HevFinalScenes{} human-verified
scenes (Table~\ref{tab:final}). They hold \HevFinalPositives{} vessels, of
which \HevFinalDarkSupport{} are dark and \HevFinalNearShoreSupport{} lie
within 2\,km of shore. F1 falls to \HevFinalFMin--\HevFinalFMax, on average
\HevFinalMeanTestGap{} below test F1. Precision stays between
\HevFinalPrecisionMin{} and \HevFinalPrecisionMax, but recall drops to
\HevFinalRecallMin--\HevFinalRecallMax: the detectors miss vessels rather
than over-detect. No cell exceeds
\HevFinalNearShoreFMax{} near-shore F1 or \HevFinalDarkRecallMax{}
dark-vessel recall. The training labels explain much of this gap: only
\HevTrainNearShorePositives{} of the \HevTrainPositives{} training vessels
(\HevTrainNearShorePct\%) lie within 2\,km of shore, against
\HevFinalNearShorePct\% of the verified vessels.

The 12-scene pretraining advantage survives this shift: the six gains over
random range from \HevDeltaFinalTenCNNOptical{} to
\HevDeltaFinalTenViTImageNet, and ViT-ImageNet at 12 scenes matches the ViT
random floor at 111 (\HevFinalFViTImageNetTen{} versus
\HevFinalFViTRandomHundred). CNN-SAR again leads optical only with few
labels (\HevFinalFCNNSarTen{} versus \HevFinalFCNNOpticalTen{} at 12 scenes,
\HevFinalFCNNSarHundred{} versus \HevFinalFCNNOpticalHundred{} at 111). Both
SAR declines from the test split recur on this set and grow larger.
\else
The 50 human-verified scenes remain sealed. Dark-vessel recall and near-shore
performance are defined only on that set, so neither result is available in
this report.
\fi
